% TOP_PROBLEM: {"id": 1439, "title": "Implicit Graph Conjecture", "category_id": 3, "category": {"id": 3, "name": "graph_theory", "display_name": "Graph Theory", "description": "Problems involving graphs, networks, and their properties.", "slug": "graph-theory", "order_index": 3, "created_at": "2026-07-31T15:26:25.671Z"}, "status": "open"}
% FIRST_POSTED: null
% ENTRY_KIND: statement_only
% Imported from: batch07.tex; UnsolvedMath ID 1439
\documentclass[11pt]{article}
\usepackage[margin=25mm]{geometry}
\usepackage{fontspec}
\setmainfont{Latin Modern Roman}
\usepackage{amsmath,amssymb,mathtools}
\usepackage{unicode-math}
\setmathfont{Latin Modern Math}
\usepackage{xurl,hyperref}
\hypersetup{colorlinks=true,urlcolor=blue}
\setcounter{secnumdepth}{0}
\setlength{\parindent}{0pt}
\setlength{\parskip}{5pt}
\setlength{\emergencystretch}{3em}
\newcommand{\sref}[2]{\hyperlink{source:#1:#2}{[S#2]}}
\newcommand{\cataloguescope}[1]{\paragraph{Recorded target.}#1}
\begin{document}
% UnsolvedMath ID 1439; source catalogue code GRAPH-052
\setcounter{section}{1438}
\section{Implicit Graph Conjecture}\label{1439}
{\small\sffamily UnsolvedMath ID 1439 · Graph Theory}
\subsection{Definitions and mathematical statement}

\paragraph{Shared notation.}$\mathbb N=\{1,2,\ldots\}$, and $\mathbb Z,\mathbb Q,\mathbb R,\mathbb C$ denote the integers, rationals, reals and complex numbers. Any other conventions are specified locally.


Let $\mathcal G$ be a family of finite simple graphs closed under graph isomorphism. It is hereditary if every induced subgraph of a member is also a member. Let $\mathcal G_n$ be its graphs with labeled vertex set $\{1,\ldots,n\}$. It has at most factorial speed when some constant $C$ satisfies
\[
 |\mathcal G_n|\le 2^{C n\log_2(n+1)}\quad(n\ge1).
\]
An implicit adjacency representation assigns to each vertex of each $G\in\mathcal G_n$ a binary label of length at most $c\lceil\log_2(n+1)\rceil$, for one constant $c$, together with a fixed decoder $D$ on pairs of binary strings such that for distinct $u,v$,
\[
 uv\in E(G)\quad\Longleftrightarrow\quad D(\ell_G(u),\ell_G(v))=1.
\]
The decoder may depend on the family but not on the individual graph. The original algorithmic convention also requires $D$ to be computable in polynomial time in the label lengths; the source's negative result permits even unrestricted decoders.
\paragraph{Statement recorded in the supplied source.}
Does every hereditary family of at most factorial speed admit an implicit adjacency representation with labels of length $O(\log n)$, in the above original algorithmic convention? The factorial growth condition is the precise standard meaning of the catalogue's “slowly-growing” family, as identified in its literature review. The same existence question with unrestricted decoders is also refuted by Source S2. \sref{1439}{2}
\subsection{Short English statement}
Must every hereditary graph family with at most factorially many labeled graphs have logarithmic-size vertex labels from which adjacency can be decoded?
\subsection{Sources}
\begin{description}
\item[\hypertarget{source:1439:1}{S1}] ULAM.AI and the UnsolvedMath Contributors, UnsolvedMath: A Curated Collection of Open Mathematics Problems, record 1439 (GRAPH-052). CC BY 4.0. \url{https://huggingface.co/datasets/ulamai/UnsolvedMath}.
\item[\hypertarget{source:1439:2}{S2}] Hamed Hatami and Pooya Hatami, The Implicit Graph Conjecture Is False, arXiv:2111.13198v2 (2021), Section 1, Conjecture 1.3 and Theorem 1.4, pages 1--3. \url{https://arxiv.org/pdf/2111.13198v2}.
\end{description}
\label{end:1439}
\end{document}
